\documentclass[11pt,a4paper]{article}

\usepackage[utf8]{inputenc}
\usepackage[T1]{fontenc}
\usepackage[a4paper,margin=2.4cm]{geometry}
\usepackage{lmodern}
\usepackage{xcolor}
\usepackage{booktabs}
\usepackage{array}
\usepackage{tabularx}
\usepackage{amsmath}
\usepackage{listings}
\usepackage{hyperref}

\definecolor{navy}{HTML}{17365D}
\definecolor{lightblue}{HTML}{EAF2F8}
\definecolor{codegray}{HTML}{F5F5F5}

\hypersetup{
  colorlinks=true,
  linkcolor=navy,
  urlcolor=navy,
  pdftitle={Análisis de viabilidad: Knowledge Graph para Dispute Transaction Support},
  pdfauthor={Factored Hackathon 2026}
}

\lstset{
  basicstyle=\ttfamily\small,
  backgroundcolor=\color{codegray},
  frame=single,
  rulecolor=\color{lightblue},
  columns=fullflexible,
  keepspaces=true,
  showstringspaces=false,
  breaklines=true,
  literate={á}{{\'a}}1 {é}{{\'e}}1 {í}{{\'i}}1 {ó}{{\'o}}1 {ú}{{\'u}}1 {ñ}{{\~n}}1 {Á}{{\'A}}1 {É}{{\'E}}1 {Í}{{\'I}}1 {Ó}{{\'O}}1 {Ú}{{\'U}}1 {Ñ}{{\~N}}1
}

\renewcommand{\arraystretch}{1.18}
\setlength{\parindent}{0pt}
\setlength{\parskip}{6pt}

\title{\textcolor{navy}{Análisis de viabilidad\\[2pt] Knowledge Graph para Dispute Transaction Support}}
\author{Factored Hackathon 2026}
\date{27 de septiembre de 2026}

\begin{document}

\maketitle

\begin{abstract}
Este informe evalúa la viabilidad de fortalecer \emph{Dispute Transaction Support} con un modelo probabilístico y un Knowledge Graph. La conclusión principal es que los datos no soportan todavía una inferencia amplia para todo el soporte general de tarjetas, pero sí justifican una primera solución especializada en \textbf{Transaction Issue / Transaction Dispute Triage}, combinando herramientas determinísticas con inferencia interpretable.
\end{abstract}

\section{Conclusión ejecutiva}

No conviene ampliar todavía la inferencia probabilística de \emph{Dispute Transaction Support} hacia todo el soporte general de tarjetas. Los datos sí alcanzan para una capa probabilística, pero el problema debe acotarse.

El mejor candidato es \textbf{Transaction Issue / Transaction Dispute Triage}, especialmente para:

\begin{itemize}
  \item transacciones rechazadas;
  \item operaciones fraudulentas o potencialmente fraudulentas;
  \item operaciones pendientes;
  \item operaciones reversadas; y
  \item casos que requieren escalamiento.
\end{itemize}

La razón es que existe bastante información operacional de tarjetas, pero poca evidencia etiquetada de soporte. Las transacciones tienen campos estructurados como \texttt{transaction\_status}, \texttt{response\_code}, \texttt{is\_fraud}, \texttt{fraud\_score}, canal, monto y ubicación. En cambio, hay poca información que represente explícitamente qué problema tenía el cliente y cómo se resolvió.

\section{Cobertura de los datos}

El conjunto contiene:

\begin{itemize}
  \item 100\,102 tarjetas de crédito y 39\,938 tarjetas de débito;
  \item más de un millón de transacciones asociadas a tarjetas de crédito;
  \item 15\,486 reclamos asociados a productos de tarjeta;
  \item 1\,223 interacciones que mencionan explícitamente IDs de tarjeta; y
  \item solo 322 transcripts asociados a esas interacciones explícitas.
\end{itemize}

Las tarjetas cuentan con contexto operacional útil: estado, saldo, límite de crédito, días de mora, fecha de última transacción, canal de apertura y vinculación con la aplicación móvil.

\subsection{Distribución transaccional de tarjetas de crédito}

La distribución observada en las transacciones de crédito es:

\begin{table}[ht]
\centering
\caption{Estados de transacciones de tarjetas de crédito}
\begin{tabular}{lrr}
\toprule
Estado & Registros & Proporción aproximada \\
\midrule
Aprobada & 1\,019\,373 & 92.0\% \\
Rechazada & 55\,650 & 5.0\% \\
Pendiente & 21\,952 & 2.0\% \\
Reversada & 11\,310 & 1.0\% \\
\bottomrule
\end{tabular}
\end{table}

Esta distribución proporciona targets transaccionales observables. Es considerablemente más sólida que intentar inferir directamente una intención general de soporte a partir de los transcripts disponibles.

\section{Viabilidad por capacidad}

\begin{table}[ht]
\centering
\small
\caption{Evaluación de capacidades para Dispute Transaction Support}
\begin{tabularx}{\textwidth}{>{\raggedright\arraybackslash}p{0.31\textwidth} X}
\toprule
Capacidad & Evaluación \\
\midrule
Consultar estado de tarjeta & \textbf{Alta}. Hay producto, estado, saldo, límite y mora. \\
Explicar transacciones rechazadas & \textbf{Alta}. Hay estado, canal, código de respuesta, monto y contexto de producto. \\
Detectar señales de fraude & \textbf{Media-alta}. Hay \texttt{is\_fraud}, \texttt{fraud\_score}, ubicación, canal y monto. \\
Identificar intención del cliente & \textbf{Baja}. Los transcripts explícitos de tarjeta son pocos y predominan en \texttt{consulta\_general}. \\
Automatizar resolución de casos & \textbf{Baja}. Faltan acciones, códigos de resolución y políticas operativas. \\
Medir satisfacción y calidad & \textbf{Baja}. La satisfacción está muy incompleta en los reclamos. \\
Crear un copiloto para agentes & \textbf{Alta}. Se pueden combinar herramientas de consulta con recomendaciones explicables. \\
Crear un agente autónomo & \textbf{Insuficiente por ahora}. Se requiere autenticación, políticas y outcomes confiables. \\
\bottomrule
\end{tabularx}
\end{table}

\section{Por qué no conviene ampliar hacia soporte general de tarjetas}

Para un modelo de soporte general se necesitarían targets como:

\begin{lstlisting}
transaction_status
response_code
fraud_score
card_status
days_past_due
customer_context
        |
        +--> support_intent
        |
        +--> resolution_action
\end{lstlisting}

Sin embargo, \texttt{support\_intent} y \texttt{resolution\_action} no están suficientemente construidos en los datos actuales. Las 1\,223 interacciones con mención explícita de tarjeta tienen 322 transcripts, y esos transcripts presentan prácticamente una sola intención dominante: \texttt{consulta\_general}.

Por ello, todavía no hay una base suficiente para entrenar de manera confiable intenciones como:

\begin{itemize}
  \item pérdida o robo;
  \item bloqueo y desbloqueo;
  \item activación;
  \item cambio de PIN;
  \item compra desconocida;
  \item retiro no reconocido;
  \item tarjeta vencida;
  \item aumento de límite; y
  \item reemplazo de tarjeta.
\end{itemize}

\section{Transaction Issue / Dispute Triage}

En el dominio transaccional la situación es más defendible estadísticamente. Las transacciones tienen volumen, variables estructuradas y estados observables. La solución podría comenzar con clases que se puedan verificar directamente:

\begin{lstlisting}
TRANSACTION_OK
DECLINED
PENDING
REVERSED
POTENTIAL_FRAUD
\end{lstlisting}

En una segunda etapa, estas clases podrían evolucionar hacia categorías de soporte:

\begin{lstlisting}
DECLINED_INSUFFICIENT_FUNDS
DECLINED_CARD_STATUS
SUSPECTED_FRAUD
UNRECOGNIZED_TRANSACTION
PENDING_TRANSACTION
REVERSED_TRANSACTION
DISPUTE_REQUIRES_ESCALATION
\end{lstlisting}

Las categorías más específicas solo deberían introducirse después de comprobar que pueden obtenerse o construirse de manera reproducible a partir de BigQuery.

\subsection{Modelo probabilístico interpretable}

Una formulación inicial podría estimar:

\[
P(\text{Fraud}\mid \text{channel},\text{amount bucket},\text{fraud score bucket},\text{response code},\text{location},\ldots)
\]

o:

\[
P(\text{Dispute Category}\mid F_1,\ldots,F_n).
\]

El Knowledge Graph podría mantener evidencia y relaciones de soporte:

\begin{lstlisting}
response_code:X  ----+
fraud_score:high ----+--> SUSPECTED_FRAUD
channel:online   ----+
amount:high      ----+
\end{lstlisting}

Cada inferencia debería conservar:

\begin{itemize}
  \item probabilidad posterior;
  \item evidencia utilizada;
  \item soporte;
  \item confianza; y
  \item lift o fuerza relativa de la asociación.
\end{itemize}

Esto mantiene la misma filosofía de un modelo Naive Bayes/Knowledge Graph, pero con millones de observaciones estructuradas en vez de un conjunto pequeño de perfiles.

\section{Arquitectura recomendada}

La propuesta debería separar claramente consultas determinísticas de inferencia probabilística:

\begin{lstlisting}
DISPUTE TRANSACTION SUPPORT
        |
        +-- Deterministic tools
        |      estado de tarjeta
        |      saldo y límite
        |      mora
        |      últimas transacciones
        |
        +-- ML inference
               |
        TRANSACTION ISSUE TRIAGE
               |
        rechazo / fraude / pending /
        reversal / dispute candidate
               |
        recommended next action
\end{lstlisting}

Las herramientas determinísticas deben ser la fuente de verdad para el estado actual de la cuenta. El modelo probabilístico debe priorizar hipótesis, detectar señales y recomendar el siguiente paso, pero no sustituir las reglas operativas del banco.

\section{Reclamos y límites para entrenamiento supervisado}

Existen 15\,486 reclamos asociados a tarjetas. No obstante:

\begin{itemize}
  \item solo 3\,534 tienen texto de resolución;
  \item solo 561 tienen satisfacción registrada;
  \item ninguno tiene \texttt{origin\_interaction\_id}; y
  \item aproximadamente 75\% continúan abiertos, en proceso o escalados.
\end{itemize}

Esto permite analizar categorías, prioridad, SLA y recurrencia, pero limita el entrenamiento de un modelo que aprenda una política completa de resolución.

\section{Recomendaciones para el proyecto}

\begin{enumerate}
  \item Presentar el producto como \textbf{Dispute Transaction Support} con herramientas determinísticas y una capa probabilística especializada en \textbf{Transaction Issue / Dispute Triage}.
  \item No presentar inicialmente el sistema como un agente que resuelve automáticamente todo tipo de reclamo.
  \item Crear una tabla normalizada \texttt{dispute\_transaction\_cases} con los campos:
\end{enumerate}

\begin{lstlisting}
customer_id
product_id
transaction_id
intent
issue_type
action_taken
resolution_code
escalated
verified_identity
outcome
\end{lstlisting}

\begin{enumerate}
  \setcounter{enumi}{3}
  \item Relacionar explícitamente conversación, tarjeta, transacción, reclamo y resultado.
  \item Incorporar una taxonomía de intenciones específica para tarjetas.
  \item Añadir una base de conocimiento con políticas de bloqueo, reemplazo, fraude, disputas, límites y tiempos de atención.
  \item Normalizar valores y tipos del contrato de datos antes del entrenamiento. Por ejemplo, el diccionario usa valores como \texttt{Credit Card}, mientras BigQuery utiliza \texttt{Tarjeta Crédito}; además, algunos tipos documentados como \texttt{INTEGER} o \texttt{VARCHAR} aparecen como \texttt{FLOAT64} o \texttt{INT64}.
\end{enumerate}

\section{Conclusión}

Para ampliar Dispute Transaction Support hacia soporte general de tarjetas no hay todavía suficiente información etiquetada para hacer una inferencia estadística tan sólida como la de un modelo de perfiles o arquetipos.

Para la parte transaccional sí existe una base adecuada: hay volumen, estados, señales de fraude, canales, montos, ubicaciones y contexto de tarjeta. Por eso, la dirección más defendible es:

\begin{quote}
\textbf{Dispute Transaction Support: herramientas determinísticas + ML probabilístico interpretable especializado en Transaction Issue / Dispute Triage.}
\end{quote}

\end{document}
